\subsection{位数4の三角境界：三つの持続的二重点と $D_4$ 衝突}
\label{subsec:quantum-triangle-even4}
$q^2=-1$ とする。奇数位数とは異なり、一般の切断中心も特異である。
この場合も差分表示・全矢は前節のものを使う。
\[
 X=x^4-qx^2-\frac18,\quad Y=y^4-qy^2-\frac18,\quad
 Z=z^4-qz^2-\frac18
\]
と置く。

\begin{theorem}[位数4pencilの全判定]
\label{thm:quantum-triangle-order4}
$V_q$ の全中心関係は
\[
 XYZ+\frac{X^2+Y^2+Z^2}{16}
 -\frac{\Omega_q^2}{2}(X+Y+Z)
 +\Omega_q^4+\frac{\Omega_q^2}{16}-\frac1{1024}=0
\]
である。$U_{q,D}$ の全中心は $\Omega_q=-D$ の商と一致する。
$u=-16X,v=-16Y,w=-16Z,k=64D^2$ で正規化すれば
\[
 F_k=uvw-u^2-v^2-w^2-2k(u+v+w)-k^2-4k+4=0
\]
となる。$k\ne0,4$ では三つの $A_1$、$k=0$ では四つの $A_1$、
$k=4$ では一つの $D_4$ を持つ。
全切断のH3条件は
\[
 R_{q,D}\text{ が強いH3}\quad\Longleftrightarrow\quad
 D^2\ne\frac1{16}
\]
である。これは $D\ne D_+,D_-$ とも一致する。
\end{theorem}
\begin{proof}
中心生成元はChebyshev中心命題で既知であり、表示関係はPBW次数12で厳密に検算できる。
$k[X,Y,Z]$ 上の $1,\Omega_q,\Omega_q^2,\Omega_q^3$ の最高記号の独立性から
他の関係はなく、$V_q$ はこの多項式部分環上自由階数64である。
固定 $D$ の表示中心は孤立特異点だけを持つ正規曲面である。
差分巡回代数へのgeneric拡大でPI次数4、中心分数体を確認し、
整性からこれが $U_{q,D}$ の全中心となる。
このgeneric拡大では $a=u_0^4$、$s=\mathcal T^4$ を使う
（差分変数 $u_0$ はここでの中心座標 $u$ と区別する）。
$x/\kappa=u_0+u_0^{-1}$ と $u_0^4=a$ から $u_0$ を、
次に $y,z$ から $\mathcal T$ を復元する。
中心分数体の拡大は、差分環の反転
$u_0\mapsto u_0^{-1}$、$\mathcal T\mapsto b_D(u_0)\mathcal T^{-1}$
による二次拡大であり、次数4巡回代数を得る。

偏微分の差は例えば $(v-u)(w+2)=0$ となる。
二座標が異なる解は $(-2,-2,2-k)$ の三つの置換であり、
これらは全 $k$ で曲面上にある。
三座標が等しい場合は $u=v=w=z$、$k=(z^2-2z)/2$ であり、
曲面式は $-(z-2)(z+2)^3/4$ となる。
従って $k=0$ で $(2,2,2)$ が追加され、$k=4$ で三点が $(-2,-2,-2)$ に合流する。
持続点のHessian行列式は $2(k-4)^2$、追加点でも非零である。
合流点の局所式は厳密に
\[
 abc-(a+b+c)^2=0
\]
であり、二次部の核上の三次式は $-ab(a+b)$ で三つの異なる因子を持つ。
従ってこれは $D_4$ 二重点である。

$U_{q,D}$ はPI次数4の対称MCM非分岐最大整環である。
その証明は先頭量子環の最大性、CY双対と余次元1Azumayaの議論による。
無限遠の中心は滑らかな三角形で、各辺のcoverは次数4、両端で完全分岐する。
前節の主radical式を $N=4$ で適用し、terminal局所模型で境界全体は有限大域次元となる。

affineの持続 $A_1$ 点には、$a=(1+q)/2$、$c=a^2-b^2$ として
\[
 x=\begin{pmatrix}a&0\\0&-a\end{pmatrix},\quad
 y=\begin{pmatrix}b&1\\c&b\end{pmatrix},\quad
 z=\begin{pmatrix}b&q\\-qc&-b\end{pmatrix}
\]
を使う。この表現では
\[
 \Omega_q=qb^2+\frac14,\qquad
 (u,v,w)=(-2,2-k,-2).
\]
任意の $D$ に対して $b^2=q(D+1/4)$ を選べる。
$c=0$ となるのは $D=1/4$ だけで、他では $x$ の相異なる固有値と
$y$ の両非対角成分の非消滅から既約である。
循環置換で持続的三点を覆う。
$D=0$ の追加 $A_1$ には零characterを使う。
従って $D\ne\pm1/4$ では全 $A_1$ がfull McKayとなり、affineの大域次元は2である。

一方、$D=\pm1/4$ の $D_4$ 被覆群は二進二面体群 $Q_8$ であり、
既約指標次元は $1,1,1,1,2$ である。
full McKayには全重複度が正であることが必要だが、
階数式 $\sum r_\chi\dim\chi=4$ の右辺より最小必要値6が大きい。
従って両値とも局所大域次元は無限となる。
AS Gorenstein・AZ・Serreと整域性の認識により主張を得る。
\end{proof}

\texttt{verify\_cubic\_quantum\_triangle\_\allowbreak order4.py} は
全中心関係・四中心生成元の交換子・特異値・Hessian・$D_4$ 局所式を厳密に検算する。
二次元表現について三関係・Casimir・全中心座標も有理式を $q^2+1$ で還元して確認する。
これは指定pencilの判定であり、偶数位数全体の判定ではない。

\subsubsection{一般偶数位数：残る一点への還元}
\begin{theorem}[偶数位数の全非端点と位数18以下の完全判定]
\label{thm:quantum-triangle-even-partial}
$q$ の正確な位数を $N=2M\ge4$ とする。
全 $D\ne D_+,D_-$ の $R_{q,D}$ は強いH3である。
$D_+$ は全偶数位数で非H3となる。
$4\le N\le18$ では $D_-$ も非H3であり、従ってこの範囲の偶数位数pencilも完全判定される。
$N\ge20$ の $D_-$ に対する結論は本定理に含めない。
\end{theorem}
\begin{proof}
$h=q^2-1$、$\Gamma_x=-h^Nc_N(x)$ cyclicとする。
Huangの中心表示\cite[Introduction, equation (5), Section 4.3]{HuangAWCentre}を使う。
規約を明記すると、同論文の量子パラメータは本稿の $\rho^{-1}$、
正規化生成元は $-x/\kappa$ cyclic、正規化Casimirは
\[
 \Pi=(q+q^{-1})-\frac{h^3}{q^2}\Omega_q
\]
である。従って切断では $\Pi=-w_D$、
$k=2+P_M(w_D)$ と置けば全中心は
\[
 F_k(\Gamma_x,\Gamma_y,\Gamma_z)=0
\]
となる。これは前定理と同じ三次式である。
中心の正規性、PI次数 $N$、対称MCM最大整環性、
無限遠三角形のterminal性は、全中心生成元のモニック性と
主radical式を用いた既証の議論で従う。
以下はこの明示的中心上でのMcKay判定である。

まず $k\ne4$ の三つの持続的 $A_1$ を覆う。
添字 $j=0,\ldots,M-1$ に対して
\[
 \begin{gathered}
 \lambda_j=-\frac{q^{j+1}+q^{-j}}h,\quad A_j=q^j-q^{1-j},\\
 c_j=\frac{-(q-1)^2b^2+q^2(q^{2j}+q^{-2j}-2)/h^2}{A_jA_{j+1}}
 \quad(1\le j<M).
 \end{gathered}
\]
と置く。$x$ は対角 $\lambda_j$、$y$ は上隣接成分1、下隣接成分 $c_j$、
両端の対角成分を $b$、他の対角成分を零とする。
$z=xy-qyx$ の上隣接成分は $q^{j+1}$、下隣接成分は $q^{-j}c_j$、
両端の対角成分は $b,-b$ である。
$M=2$ でも両端を異なる位置としてこの定義を使う。
三関係は内部再帰式と最初・最後の対角式で厳密に成立し、
\[
 \Omega_q=\frac{q-1}{q+1}b^2+\frac q{(q-1)^3(q+1)},\qquad
 w_D=2+\frac{(q-1)^4(q+1)^2}{q^2}b^2
\]
となる。従って全 $D$ を覆う。
$A_j,A_{j+1}$ の非消滅は $N$ が偶数で指数が奇数であることによる。
$\lambda_j$ は相異なり、$c_j=0$ は
$w_D=q^{2j}+q^{-2j}$ と同値、すなわち $k=4$ のときだけ起こる。
従って $k\ne4$ では次元 $M$ の既約表現である。
その $\Gamma_x=-2$ は対角成分から直接分かるので、
$k=0$ の追加点 $(2,2,2)$ には行かず、持続点の一つへ行く。
循環置換で三点全てを覆い、各点はfull $A_1$ McKayとなる。

$k=0$ の端点以外の値では、前節の古典型表示を
$1\le r<M$、$w_D=-(q^r+q^{-r})$ で使う。
その分母・隣接成分は非零、$x$ の固有値も相異なることを同じ指数範囲で確認できる。
この表示は追加点へ行く。
実際、持続点へ行けば既にある次元 $M$ のsimpleと異なる次元 $r<M$ のsimpleを持つが、
$A_1$ の重複度は二つ以下で和は $2M$ なので矛盾する。
従って追加点もfull McKayである。

$k=4$ の端点以外では一つの $D_4$ 点を持つ。
$w_D=q^{2d}+q^{-2d}$、$1\le d<M/2$ と書ける。
前節の非古典型次元 $d$ 表示はこの範囲でも既約であり、
二符号変更による四つのtraceベクトルは
$(-d_0,d_0,d_0)$ の二符号変更で、$d_0\ne0$ なので全て相異なる。
一方、古典型次元 $r=M-2d>0$ の表示も同じ $w_D$ を持ち、traceは三生成元とも零である。
従って一つの $D_4$ 中心点に五つの相異なるsimpleを得る。
これは被覆群 $Q_8$ の全五指標に一致するのでfull McKayである。
これで全非端点の許容性が従う。

残る端点を調べる。$D_+$ は $w_D=-2$ に対応する。
$M$ が奇数なら $k=0$ の追加 $A_1$ で三次元中心は正則であり、
$P_M'(-2)\ne0$ と $\partial F_k/\partial k\ne0$ からAzumayaの特異中心商となる。
従って非H3である。
$M$ が偶数なら $k=4$ となり、非古典型 $d=M/2$ の四つの既約表示が存在する。
全てのsimple次元が $d$ なので、これら四つが $D_4$ の外側四指標なら
既にrank寄与 $4d=2M=N$ を使い切る。
一つが二次元指標なら寄与は少なくとも $5d>N$ となる。
いずれも全五指標正重複度は不可能なので非H3である。

$D_-$ は $w_D=2,k=4$ であり、古典型 $r=M$ の表示が既約のまま存在する。
また、この中心切断には一次元表現は一つもない。
全characterは零点と四つの符号点で、その $w_D$ は
$-(q+q^{-1})$ と $q^2+q^{-2}$ であり、いずれも2と異なるからである。
従ってfull $D_4$ なら全重複度は2以上である。
次元 $M$ のsimpleが二次元指標に属せばそれだけでrank $2M$ を使い切り、
一次元指標に属せば最小総rankは
\[
 M+3\cdot2+2\cdot2=M+10
\]
となる。$M\le9$ では $M+10>2M$ なのでfull McKayは不可能である。
これが位数18以下の完全判定を与える。
$M\ge10$ ではこの下界だけで結論できないため、$D_-$ の判定を保留する。
\end{proof}

\texttt{verify\_cubic\_quantum\_triangle\_\allowbreak even\_representations.py} は、
$j$ を形式変数 $J=q^j$ に置き換えた全内部再帰、両端の式、Casimirと
既約性を失う係数の消滅方程式を全有理式で厳密に検算する。
位数による指数の非消滅とMcKay指標数・rank下界は本文の証明であり、標本計算による外挿ではない。
